
\documentclass{article} % For LaTeX2e
\usepackage{iclr2027_conference,times}

% Optional math commands from https://github.com/goodfeli/dlbook_notation.
\input{math_commands.tex}

\usepackage{hyperref}
\usepackage{url}

\usepackage[utf8]{inputenc}
\usepackage[T1]{fontenc}
\usepackage{booktabs}
\usepackage{amsfonts}
\usepackage{amsmath}
\usepackage{amssymb}
\usepackage{amsthm}
\usepackage{bm}
\usepackage{microtype}
\usepackage{xcolor}
\usepackage{colortbl}
\usepackage{graphicx}
\usepackage{algorithm}
\usepackage{algorithmic}
\usepackage{multirow}
\usepackage{makecell}
\usepackage{placeins}

\newcommand{\softmin}{\operatorname*{softmin}}
\newcommand{\subfloat}[2][]{%
  \begin{minipage}[t]{0.30\textwidth}
    \centering
    #2\\[0.4em]
    \footnotesize #1
  \end{minipage}%
}
\newcommand{\missingfigureplaceholder}{%
  \setlength{\fboxsep}{0pt}%
  \fbox{%
    \begin{minipage}[c][0.18\textheight][c]{0.28\textwidth}
      \centering
      \small Figure placeholder\\[0.4em]
      \footnotesize Figure unavailable
    \end{minipage}%
  }%
}
\newcommand{\safeincludegraphics}[2][]{%
  \IfFileExists{#2}{\includegraphics[#1]{#2}}{\missingfigureplaceholder}%
}
\newtheorem{proposition}{Proposition}[section]

\title{CARAT: Class-Aware Robust Aggregation for Model-Poisoning-Resilient Federated Learning}

% Authors must not appear in the submitted version. They should be hidden
% as long as the \iclrfinalcopy macro remains commented out below.
% Non-anonymous submissions will be rejected without review.

\author{Antiquus S.~Hippocampus, Natalia Cerebro \& Amelie P. Amygdale \thanks{ Use footnote for providing further information
about author (webpage, alternative address)---\emph{not} for acknowledging
funding agencies.  Funding acknowledgements go at the end of the paper.} \\
Department of Computer Science\\
Cranberry-Lemon University\\
Pittsburgh, PA 15213, USA \\
\texttt{\{hippo,brain,jen\}@cs.cranberry-lemon.edu} \\
\And
Ji Q. Ren \& Yevgeny LeNet \\
Department of Computational Neuroscience \\
University of the Witwatersrand \\
Joburg, South Africa \\
\texttt{\{robot,net\}@wits.ac.za} \\
\AND
Coauthor \\
Affiliation \\
Address \\
\texttt{email}
}

% The \author macro works with any number of authors. There are two commands
% used to separate the names and addresses of multiple authors: \And and \AND.
%
% Using \And between authors leaves it to \LaTeX{} to determine where to break
% the lines. Using \AND forces a linebreak at that point. So, if \LaTeX{}
% puts 3 of 4 authors names on the first line, and the last on the second
% line, try using \AND instead of \And before the third author name.

\newcommand{\fix}{\marginpar{FIX}}
\newcommand{\new}{\marginpar{NEW}}

% \iclrfinalcopy % Uncomment for camera-ready version, but NOT for submission.
\begin{document}


\maketitle





\begin{abstract}
Federated learning is vulnerable to model-poisoning attacks because the server must aggregate client updates without inspecting client data. Existing robust aggregation rules often rely on update geometry, coordinate-wise robust statistics, or a single trusted reference direction, but these signals can confuse malicious behavior with benign client heterogeneity under non-i.i.d.\ data. We propose \textbf{CARAT}, a class-aware robust aggregation rule that uses labeled server-side examples to sample hidden, class-balanced tasks and score candidate updates by task-level loss reduction. CARAT combines these certificates with robust clipping, common-radius evaluation, coordinate-rank anomaly penalties, and capped-simplex weight optimization. We evaluate CARAT over five seeds on CIFAR-100/ResNet18 with $20\%$ malicious clients, three non-i.i.d.\ regimes, and five attacks. For $\alpha\in\{1,0.5\}$, CARAT has average rank $2.5$ under four conventional attacks and remains within $0.1$--$1.2$ percentage points of the strongest mean accuracy. Under Mimic, it ranks second among defenses with recorded results. At $\alpha=0.1$, however, CARAT becomes highly seed-sensitive under FangAttack and MinSum, exposing a clear robustness boundary. A component study further finds only $0.06$--$0.28$ percentage-point separation between the complete method and reduced variants. The evidence therefore supports conditional competitiveness, not universal superiority or large independent gains from every component.
\end{abstract}


\section{Introduction}
\label{sec:introduction}

Federated learning (FL) enables many clients to train a shared model without centralizing their private data \citep{mcmahan2017fedavg,bonawitz2017secureagg,kairouz2021federated}. This distributed training paradigm is attractive for privacy- and governance-sensitive applications, but it also makes aggregation a security-critical operation. At each round, the server must update the global model from client messages whose underlying data and training process cannot be directly inspected. As a result, the aggregation rule becomes the main line of defense when some clients submit harmful updates.

Model-poisoning attacks exploit this trust bottleneck by replacing benign client updates with crafted messages. These attacks can be targeted, such as backdoor and model-replacement attacks \citep{bagdasaryan2020backdoor,xie2020dba,wang2020tails}, or untargeted, where the goal is to degrade clean test performance and destabilize collaborative training \citep{baruch2019alie,fang2020local,shejwalkar2021manipulating}. This paper focuses on the untargeted setting. The core challenge is not only that malicious updates can be arbitrary, but also that benign updates are naturally diverse under non-i.i.d.\ client data. A geometrically unusual update may therefore be either an attack or a legitimate consequence of client heterogeneity \citep{zhao2018noniid,li2020fedprox,karimireddy2020scaffold,zhu2021noniidsurvey}.

Existing Byzantine-robust defenses address this threat through geometric selection, coordinate-wise robust statistics, or auxiliary validation, but their trust signals can be too coarse for heterogeneous FL. Krum and MultiKrum \citep{blanchard2017krum} select centrally located updates, Bulyan \citep{elmhamdi2018bulyan} combines selection with a second robust aggregation stage, coordinate-wise median \citep{yin2018median} suppresses marginal outliers, and RFA \citep{pillutla2022rfa} uses a geometric-median objective. These rules can confuse benign dispersion with malicious deviation when client data are non-i.i.d. Auxiliary-information defenses partly address this limitation, but FLTrust \citep{cao2021fltrust} compresses trusted data into a single reference direction and FLDetector \citep{zhang2022fldetector} relies on historical behavior prediction. Such signals may miss attacks whose harm is concentrated on a subset of classes or semantic regions. The resulting open problem is how to validate client updates by their task-level usefulness rather than only by their location in update space.

We propose \textbf{CARAT}, a class-aware robust aggregation rule that validates updates through hidden, class-balanced tasks sampled from a small labeled server reference set. The key idea is that a benign non-i.i.d.\ update need not be geometrically central, but it should provide useful evidence on semantic validation tasks that are hidden from clients. CARAT first places candidate updates in a common scoring geometry through robust clipping and common-radius evaluation. It then measures task-level loss reduction to form hidden-task certificates, penalizes coordinate-rank anomalies, and solves a capped-simplex weight optimization problem with temporal regularization. The final update is computed from clipped client updates, so the server can preserve useful benign diversity while limiting updates that are semantically harmful or structurally extreme. This design explicitly targets settings where a modest labeled server reference set is available; it is not intended to be a fully data-free defense.

We evaluate CARAT on CIFAR-100~\citep{krizhevsky2009learning} with ResNet18~\citep{he2016deep} under five untargeted model-poisoning attacks, three Dirichlet non-i.i.d.\ regimes, and a $20\%$ malicious-client ratio. In the completed five-seed sweep, CARAT is competitive for $\alpha\in\{1,0.5\}$: under the four conventional attacks it has average rank $2.5$ and remains within $0.1$--$1.2$ percentage points of the strongest row. It also ranks second among defenses with Mimic results. This behavior does not extend uniformly to $\alpha=0.1$, where FangAttack causes four of five seeds to finish at chance accuracy and MinSum causes one similar collapse. The component study further shows only modest separation among variants. We therefore make a bounded claim: task-conditioned validation is a viable reference-assisted signal in moderate heterogeneity, while the current design is unstable under some attacks in the most heterogeneous tested setting. Our main contributions are as follows.

\begin{itemize}
\item We introduce hidden-task certificates for robust FL aggregation, using class-balanced server-side validation tasks to evaluate whether each candidate update improves semantic slices of the learning problem.
\item We develop the CARAT aggregation objective, which combines hidden-task utility, coordinate-rank anomaly penalties, capped per-client influence, and temporal regularization into a single robust weighting rule.
\item We provide an auditable five-seed CIFAR-100/ResNet18 evaluation against public robust-aggregation baselines under five attacks and three non-i.i.d.\ partitions, with explicit reporting of both competitive regimes and observed failure cases.
\end{itemize}

\section{Related Work}
\label{sec:related-work}

\paragraph{Model-poisoning attacks.}
Model-poisoning attacks directly manipulate client model updates before aggregation; in this paper we focus on untargeted availability attacks that prevent convergence or degrade clean test performance. Early Byzantine and distributed-learning work studied large-noise attacks in Byzantine-tolerant gradient descent \citep{blanchard2017krum}, sign-flipping attacks in asynchronous Byzantine SGD \citep{damaskinos2018asynchronous}, and local model manipulation in federated learning \citep{bhagoji2019adversarial}, establishing that robust aggregation must tolerate arbitrary client messages rather than only stochastic training noise. Later attacks became substantially more adaptive. ALIE \citep{baruch2019alie} exploits benign gradient variance to remain statistically plausible, IPM \citep{xie2020ipm} reverses progress through inner-product manipulation, FangAttack \citep{fang2020local} optimizes local model poisoning against specific robust aggregators, and MinMax/MinSum \citep{shejwalkar2021manipulating} search for malicious updates that satisfy distance- or sum-based stealth constraints while damaging the global model. Mimic-style attacks \citep{karimireddy2022bucketing} further show that heterogeneity itself can be weaponized by making poisoned updates resemble legitimate client drift. Targeted backdoor and model-replacement attacks form an important adjacent threat family, including model replacement \citep{bagdasaryan2020backdoor}, distributed backdoor attacks \citep{xie2020dba}, and tail-class backdoors \citep{wang2020tails}, but CARAT's empirical claim is deliberately limited to untargeted model poisoning.

\paragraph{Defenses against model poisoning.}
Defenses against model poisoning differ mainly in the trust signal used for aggregation. Classical Byzantine-robust rules rely on geometry or coordinate-wise statistics: Krum and MultiKrum \citep{blanchard2017krum} select centrally located updates, Bulyan \citep{elmhamdi2018bulyan} combines selection with a second robust aggregation stage, coordinate-wise median \citep{yin2018median} suppresses marginal outliers, and RFA \citep{pillutla2022rfa} uses a geometric-median objective. Subsequent defenses add suspicion scoring through Zeno \citep{xie2019zeno}, redundant gradients through DRACO \citep{chen2018draco}, regularized aggregation through RSA \citep{li2019rsa}, divide-and-conquer filtering through DnC \citep{shejwalkar2021manipulating}, history-aware clipping through CenteredClipping \citep{karimireddy2021history}, sign-pattern filtering through SignGuard \citep{xu2022signguard}, latent-space screening through FLARE \citep{wang2022flare}, and layer-adaptive sparsification through LASA \citep{xu2025lasa}. Because non-i.i.d.\ data can make benign updates geometrically dispersed, related work also studies optimization and robustness under heterogeneity through FedProx \citep{li2020fedprox}, SCAFFOLD \citep{karimireddy2020scaffold}, resampling \citep{he2021resampling}, bucketing \citep{karimireddy2022bucketing}, local-SGD analysis \citep{data2021localsgd}, gradient splitting \citep{liu2023gas}, nearest-neighbor mixing \citep{allouah2023fixing}, label-skew-aware aggregation \citep{bao2024boba}, and entropy-guided public-data aggregation \citep{huang2024sdea}; non-i.i.d.\ FL surveys provide the broader context \citep{zhu2021noniidsurvey}. Auxiliary-information defenses use trusted data, historical behavior, or diagnostic signals rather than only instantaneous update geometry: AUROR \citep{shen2016auror} screens updates in collaborative learning, FLTrust \citep{cao2021fltrust} bootstraps trust from a server reference update, FoolsGold \citep{fung2020foolsgold} penalizes repeated sybil-like similarity, and FLDetector \citep{zhang2022fldetector} predicts suspicious client behavior from history. CARAT follows this auxiliary-information direction but changes the validation principle: instead of trusting one anchor direction, one geometry statistic, or one historical detector, it asks whether each update remains useful across multiple hidden, class-balanced semantic tasks under a matched server-reference budget.

\section{Preliminaries and Problem Statement}
\label{sec:problem}

\subsection{Federated Learning and Threat Model}

We study synchronous FL in which the server must update a global model from client messages without inspecting client data. Let $\mathcal{N}=\{1,\dots,N\}$ be the client population. Client $i$ holds local data $\mathcal{D}_i$ and objective $F_i(w)$, giving the standard weighted objective
\begin{equation}
F(w)=\sum_{i=1}^N p_i F_i(w),
\qquad
p_i=\frac{|\mathcal{D}_i|}{\sum_{j=1}^N|\mathcal{D}_j|}.
\label{eq:fl-objective}
\end{equation}
At communication round $r$, the server broadcasts $w_r$ to selected clients $\mathcal{S}_r$. Each selected client initializes $w_{i,r}^{0}=w_r$ and performs $K$ local stochastic-gradient steps,
\begin{equation}
w_{i,r}^{k+1}=w_{i,r}^{k}-\eta_r \nabla F_i(w_{i,r}^{k};\xi_{i,r}^{k}),
\qquad k=0,\dots,K-1.
\label{eq:local-update}
\end{equation}
The returned model message is converted to an update vector
\begin{equation}
\Delta_i^r=w_{i,r}^{K}-w_r,
\label{eq:delta-def}
\end{equation}
and a server-side aggregation rule maps the received updates to the next global model,
\begin{equation}
w_{r+1}=w_r+\mathcal{A}\big(\{\Delta_i^r:i\in\mathcal{S}_r\}\big).
\label{eq:generic-aggregation}
\end{equation}
FedAvg is recovered when $\mathcal{A}$ is the sample-size-weighted mean.

The adversary controls a subset $\mathcal{M}_r\subseteq\mathcal{S}_r$ of selected clients and may replace their benign updates with arbitrary crafted messages $\widetilde{\Delta}_i^r$. We focus on untargeted model poisoning: malicious updates aim to degrade clean test accuracy while remaining plausible enough to evade robust aggregation. This setting is difficult under non-i.i.d.\ data because benign clients can naturally differ in update direction, norm, sparsity, and class emphasis. A robust aggregator must therefore reduce harmful influence without treating every geometrically atypical update as malicious.

The adversary is assumed to know the training protocol, defense family, and public hyperparameters, but not the realized hidden validation tasks sampled by the server in a given round. This defines CARAT's reference-assisted security boundary: the labeled reference set and task sampler are server-side resources, not public client inputs. Attackers that directly observe the reference examples or optimize against the exact hidden tasks are outside the empirical claim set; CARAT still limits accepted influence through clipping, rank penalties, and capped weights.

\subsection{Reference-Assisted Robust Aggregation}

CARAT assumes that the server maintains a small labeled reference set $\mathcal{D}^{\mathrm{ref}}$ covering the task label space. From this set it constructs a probe pool $\mathcal{P}$ and samples hidden, class-balanced validation tasks during aggregation. This is a stronger setting than fully data-free robust aggregation, and we treat it as part of the threat model. The distinction from single-anchor trusted-data methods is that CARAT does not compress $\mathcal{D}^{\mathrm{ref}}$ into one trusted update direction; instead, it uses the reference data to instantiate multiple semantic checks that are hidden from clients.

Given the round-level update set $\{\Delta_i^r\}_{i\in\mathcal{S}_r}$ and reference data $\mathcal{D}^{\mathrm{ref}}$, the server constructs nonnegative aggregation weights and an aggregate update that preserve useful benign diversity while limiting poisoned updates. This goal imposes four design requirements. First, update scoring should be scale-controlled so that large-norm messages cannot receive high trust solely by magnitude. Second, trust should depend on task-conditioned utility, not only on distance to other clients. Third, the method should still penalize structurally extreme updates, because a harmful update may look useful on part of the validation signal. Finally, the aggregation weights should have bounded per-client influence and temporal stability across rounds. CARAT implements these requirements through common-geometry scoring, hidden-task certificates, coordinate-rank penalties, and capped-simplex optimization.

\section{CARAT}
\label{sec:method}

\subsection{Overview}

CARAT turns the requirements above into the three-stage server-side pipeline shown in Figure~\ref{fig:carat-overview}. First, the server constructs hidden class-balanced tasks from a labeled reference set, making the validation queries semantic and unobserved by clients. Second, CARAT evaluates each candidate update on these hidden tasks after placing updates in a common scoring geometry, producing a certificate matrix of task-level loss reductions. Third, CARAT learns bounded aggregation weights by optimizing worst-task utility under rank-anomaly penalties, temporal stability, and per-client influence caps. The final server update is a weighted average of clipped client updates, not of the normalized evaluation copies.

\begin{figure*}[!htbp]
\centering
\safeincludegraphics[width=\textwidth,trim=12 75 18 25,clip]{figures/Carat_design.png}
\caption{Overview of CARAT. Under non-i.i.d.\ data, benign diversity can resemble poisoning. CARAT addresses this ambiguity through three server-side stages: (A) hidden class-balanced task construction from a labeled reference set, (B) task-conditioned validation that forms a certificate matrix of loss reductions, and (C) bounded weight learning with rank-anomaly penalties, temporal stability, and per-client caps. The server updates the model using a weighted average of clipped client updates.}
\label{fig:carat-overview}
\end{figure*}

Appendix~\ref{sec:appendix-component-rationale} summarizes how each component maps to a concrete failure mode in non-i.i.d.\ robust aggregation.

\subsection{Hidden Tasks and Task-Conditioned Validation}

Within a round, CARAT first prepares client updates for task-conditioned validation. We suppress the round superscript in this section. Before semantic evaluation, the server clips each update so that a small number of large-magnitude messages cannot dominate the comparison stage. Let $\|\Delta_i\|_2$ denote the norm of client $i$'s update. In our implementation, the clipping threshold $\tau_r$ is estimated by the MAD-based rule
\begin{equation}
\tau_r = \operatorname{median}(\{\|\Delta_i\|_2\}) +
k \cdot 1.4826 \cdot \operatorname{MAD}(\{\|\Delta_i\|_2\}),
\label{eq:clip}
\end{equation}
and the clipped update is
\begin{equation}
\bar{\Delta}_i = \min\left(1, \frac{\tau_r}{\|\Delta_i\|_2 + \varepsilon}\right)\Delta_i.
\end{equation}
CARAT then rescales each clipped update to a common evaluation radius $q_r$,
\begin{equation}
\tilde{\Delta}_i = q_r \frac{\bar{\Delta}_i}{\|\bar{\Delta}_i\|_2 + \varepsilon},
\label{eq:normalize}
\end{equation}
where $q_r$ is chosen as a quantile of the positive clipped norms. The normalized vector $\tilde{\Delta}_i$ is used only to evaluate candidate utility; the final aggregation step uses $\bar{\Delta}_i$. This separation makes certificate scoring scale-neutral without changing the conservative clipped update applied to the global model.

CARAT then constructs hidden semantic validation tasks from the server reference data. It samples $T$ hidden tasks from the probe pool $\mathcal{P}$, with each task constructed to be approximately class-balanced. This class-balanced construction is the source of CARAT's class awareness: a malicious update may remain globally plausible while degrading only a subset of classes, whereas repeated balanced tasks make localized failures more visible to the server. For task $t$, let $\ell_t(w)$ denote the average cross-entropy loss of model $w$ on that task. CARAT defines a \emph{hidden-task certificate}
\begin{equation}
c_{i,t} = \ell_t(w_r) - \ell_t(w_r + \tilde{\Delta}_i).
\label{eq:cert}
\end{equation}
Here, ``certificate'' denotes an internal task-level utility score, not a formal certified-robustness guarantee. Positive values indicate that client $i$ improves task $t$ relative to the current global model.

Stacking these scores yields a certificate matrix $C \in \R^{n \times T}$. Because tasks may have different intrinsic difficulty and baseline loss scale, CARAT robustly standardizes each column of $C$ before optimization. This preserves within-task ordering while making scores comparable across hidden slices. The matrix gives the server a richer signal than client-to-client similarity alone: an update receives high weight when it remains broadly useful across hidden tasks, not merely because it lies near the geometric center of the client population.

Task utility alone is insufficient, because a malicious client may look useful on part of the semantic signal while remaining structurally extreme in parameter space. CARAT therefore augments hidden-task certificates with a coordinate-rank anomaly score. For a sampled coordinate set $S \subseteq \{1,\dots,d\}$, it ranks clients coordinate-wise and measures the average normalized distance from the center rank:
\begin{equation}
\gamma_i^{(S)} = \frac{1}{|S|}\sum_{j \in S}
\left|
\frac{\operatorname{rank}_{i,j}}{n-1} - \frac{1}{2}
\right|.
\label{eq:rank}
\end{equation}
CARAT repeats this computation over several random coordinate subsets and aggregates the resulting values robustly. After robust standardization, only positive deviations are retained as the rank-penalty vector $\gamma\in\R^n$. This signal is intentionally asymmetric: suspicious extremeness is penalized, whereas central ranks are not explicitly rewarded. Together, the certificate matrix and rank penalty separate two notions that are often conflated in robust aggregation: whether an update is useful and whether it is structurally atypical.

\subsection{Bounded Weight Learning}

CARAT maps task-conditioned validation signals to aggregation weights through a bounded optimization problem on the capped simplex
\begin{equation}
\mathcal{W}(\hat{\rho}) =
\left\{
\alpha \in \R^n :
\alpha_i \ge 0,\;
\sum_{i=1}^n \alpha_i = 1,\;
\alpha_i \le \frac{1}{(1-\hat{\rho})n}
\right\},
\label{eq:capped-simplex}
\end{equation}
where $\hat{\rho}$ is an estimate of the attacker fraction. The cap prevents any single client from dominating the aggregate and encodes a conservative trust budget at the weight level: even when a client appears useful, it cannot absorb arbitrary influence within one round.

Let $\alpha^{\mathrm{prev}}$ be the previous round's weights. CARAT uses the quadratic temporal-variation penalty $\|\alpha-\alpha^{\mathrm{prev}}\|_2^2$ and solves
\begin{equation}
\max_{\alpha \in \mathcal{W}(\hat{\rho})}
\;\softmin_{\beta}(C^\top \alpha)
- \lambda_r \gamma^\top \alpha
- \lambda_p \|\alpha - \alpha^{\mathrm{prev}}\|_2^2,
\label{eq:objective}
\end{equation}
where
\begin{equation}
\softmin_{\beta}(s_1,\dots,s_T)
= -\frac{1}{\beta}\log\left(\frac{1}{T}\sum_{t=1}^T e^{-\beta s_t}\right).
\label{eq:softmin}
\end{equation}
The first term is a smooth approximation to worst-task utility: large $\beta$ places more emphasis on the weakest hidden tasks, while finite $\beta$ preserves optimization stability. The second term penalizes structurally extreme updates, and the third term regularizes the solution toward the previous round to reduce unnecessary oscillation. Eq.~\eqref{eq:objective} implements a simple decision rule: a candidate should retain influence only if it remains useful across hidden tasks and does not appear as a consistent rank-space outlier.

After optimization, CARAT aggregates the clipped updates:
\begin{equation}
\Delta_{\mathrm{agg}} = \sum_{i=1}^n \alpha_i \bar{\Delta}_i.
\label{eq:aggregate}
\end{equation}
This aggregate is the update applied to the global model. It inherits the robustness benefits of clipping while using weights learned from the richer evaluation geometry defined above. In practice, the main additional computation comes from evaluating $n$ candidate updates over $T$ hidden tasks and computing rank penalties on coordinate subsamples. Probe-model forward passes dominate, while the ranking and optimization stages remain tunable through explicit budget parameters.

A compact round-level algorithm summary is provided in Appendix~\ref{sec:appendix-algorithm}. Structural properties of feasible weighting, clipped aggregation, malicious-weight influence, hidden-task evaluation geometry, objective behavior, and per-round complexity are deferred to Appendix~\ref{sec:appendix-theory}.

\section{Experiments}
\label{sec:experiments}

We conduct experiments on CIFAR-100~\citep{krizhevsky2009learning} with ResNet18~\citep{he2016deep}. The class-rich label space provides a direct test of CARAT's class-balanced hidden tasks. We partition the training set among $20$ clients using a Dirichlet distribution with $\alpha\in\{1.0,0.5,0.1\}$, where smaller concentrations produce greater client heterogeneity. All clients participate in every round. Each run lasts $200$ communication rounds, and clients perform four local epochs of SGD per round with learning rate $0.05$, momentum $0.9$, weight decay $5\times10^{-4}$, and batch size $64$.

We evaluate five untargeted model-poisoning attacks: ALIE~\citep{baruch2019alie}, FangAttack~\citep{fang2020local}, the MinMax and MinSum attacks~\citep{shejwalkar2021manipulating}, and Mimic~\citep{karimireddy2022bucketing}. The first four construct malicious updates to evade statistical or distance-based screening; Mimic instead copies the update of a fixed benign client to create a benign-looking aggregation bias. Malicious clients constitute $10\%$ of the federation, corresponding to two of the $20$ participants. We compare CARAT against unprotected Mean aggregation, NormClipping, MultiKrum~\citep{blanchard2017krum}, FLTrust~\citep{cao2021fltrust}, and FLDetector~\citep{zhang2022fldetector}.

CARAT uses eight hidden tasks, four probe examples per class and task, a $512$-example probe pool, and a probe batch size of $128$. We set the probe-radius quantile to $0.5$, the MAD clipping factor to $2.5$, the certificate temperature to $12$, the rank and temporal-prior weights to $0.20$ and $0.05$, and the attacker-fraction prior to the true experimental value $\hat{\rho}=0.1$. Appendix~\ref{sec:appendix-experiments} provides the complete configuration and result provenance.

All reported cells use five independent runs with seeds $42$--$46$. Table~\ref{tab:cifar100-main-acc} reports final-round accuracy as mean and sample standard deviation. The available archive does not contain Mimic runs for unprotected Mean, so that cell is marked unavailable rather than inferred. To keep the ranking comparable across all six methods, average rank is therefore computed over ALIE, FangAttack, MinMax, and MinSum only. CARAT and FLTrust also use different server-reference budgets: CARAT draws a $512$-example pool and FLTrust uses $100$ examples. In both cases, these examples come from the same public training set used to form client partitions and are not removed from client data. We interpret their rows as the performance of the recorded method--reference configurations, not as an isolated comparison of aggregation rules.

\begin{table}[H]
  \definecolor{caratrow}{gray}{0.88}
  \renewcommand{\arraystretch}{1.03}
  \caption{Final test accuracy (\%) under five model-poisoning attacks with $10\%$ malicious clients on CIFAR-100/ResNet18. Values are mean$\pm$sample standard deviation over five seeds at round 200; higher is better. Average rank is computed over the four conventional attacks because the Mimic--Mean cell is unavailable. Bold values are the highest recorded attack means.}
  \label{tab:cifar100-main-acc}
  \setlength{\tabcolsep}{3.4pt}
  \centering\small
  \resizebox{\textwidth}{!}{%
  \begin{tabular}{c c l c c c c c c}
  \toprule
  \multirow{2}{*}{\makecell[c]{Dataset\\Model}} &
  \multirow{2}{*}{\makecell[c]{Non-IID\\Setting}} &
  \multirow{2}{*}{\makecell[c]{Aggregation\\Rule}} &
  \multicolumn{5}{c}{Attacks} &
  \multirow{2}{*}{\makecell[c]{Avg. Rank\\(4 attacks)}} \\
  \cmidrule(lr){4-8}
  & & & ALIE & FangAttack & MinMax & MinSum & Mimic & \\
  \midrule
  \multirow{6}{*}{\makecell[c]{CIFAR-100\\ResNet18}}
  & \multirow{6}{*}{$\alpha=1.0$}
  & Mean & 51.43$\pm$0.56 & \textbf{51.00$\pm$0.21} & 51.30$\pm$0.31 & 51.51$\pm$0.58 & -- & \textbf{2.00} \\
  & & NormClipping & 25.89$\pm$0.50 & 24.29$\pm$0.71 & 25.36$\pm$0.40 & 25.14$\pm$0.60 & 27.35$\pm$0.64 & 6.00 \\
  & & MultiKrum & 43.72$\pm$0.52 & 37.55$\pm$1.27 & 37.40$\pm$0.64 & 37.71$\pm$1.19 & 32.79$\pm$2.49 & 3.75 \\
  & & FLTrust & 33.80$\pm$0.60 & 33.93$\pm$0.65 & 34.08$\pm$0.58 & 34.13$\pm$0.50 & 33.70$\pm$0.52 & 4.75 \\
  & & FLDetector & \textbf{51.85$\pm$0.65} & 31.55$\pm$0.57 & \textbf{51.70$\pm$0.56} & \textbf{51.68$\pm$0.58} & \textbf{51.73$\pm$0.60} & \textbf{2.00} \\
  & & \cellcolor{caratrow}\textbf{CARAT ($T=8$)} & \cellcolor{caratrow}51.40$\pm$0.84 & \cellcolor{caratrow}50.97$\pm$0.51 & \cellcolor{caratrow}51.52$\pm$0.69 & \cellcolor{caratrow}51.34$\pm$0.76 & \cellcolor{caratrow}51.27$\pm$0.34 & \cellcolor{caratrow}2.50 \\
  \cmidrule(lr){1-9}
  \multirow{6}{*}{\makecell[c]{CIFAR-100\\ResNet18}}
  & \multirow{6}{*}{$\alpha=0.5$}
  & Mean & \textbf{50.45$\pm$0.86} & 49.65$\pm$0.42 & \textbf{50.57$\pm$0.69} & 50.30$\pm$0.81 & -- & \textbf{1.50} \\
  & & NormClipping & 24.68$\pm$0.71 & 23.26$\pm$0.26 & 24.08$\pm$0.35 & 24.12$\pm$0.58 & 26.25$\pm$0.72 & 6.00 \\
  & & MultiKrum & 40.40$\pm$1.28 & 33.47$\pm$2.35 & 33.36$\pm$1.81 & 33.74$\pm$1.80 & 29.20$\pm$1.14 & 3.75 \\
  & & FLTrust & 32.62$\pm$0.71 & 32.66$\pm$0.72 & 32.79$\pm$0.59 & 32.84$\pm$0.79 & 32.28$\pm$0.71 & 4.75 \\
  & & FLDetector & 50.44$\pm$0.80 & 29.74$\pm$1.98 & 50.37$\pm$0.59 & \textbf{50.50$\pm$0.95} & \textbf{50.83$\pm$0.94} & 2.50 \\
  & & \cellcolor{caratrow}\textbf{CARAT ($T=8$)} & \cellcolor{caratrow}50.09$\pm$1.00 & \cellcolor{caratrow}\textbf{49.70$\pm$0.60} & \cellcolor{caratrow}50.05$\pm$0.46 & \cellcolor{caratrow}49.96$\pm$0.60 & \cellcolor{caratrow}49.91$\pm$0.68 & \cellcolor{caratrow}2.50 \\
  \cmidrule(lr){1-9}
  \multirow{6}{*}{\makecell[c]{CIFAR-100\\ResNet18}}
  & \multirow{6}{*}{$\alpha=0.1$}
  & Mean & \textbf{44.62$\pm$0.73} & \textbf{43.92$\pm$0.47} & 44.56$\pm$0.45 & 44.32$\pm$0.42 & -- & \textbf{2.00} \\
  & & NormClipping & 20.57$\pm$0.95 & 19.55$\pm$0.90 & 19.94$\pm$0.86 & 19.80$\pm$0.90 & 22.18$\pm$0.87 & 6.00 \\
  & & MultiKrum & 31.22$\pm$1.07 & 23.18$\pm$1.76 & 23.06$\pm$1.67 & 23.38$\pm$1.95 & 19.74$\pm$1.43 & 4.75 \\
  & & FLTrust & 28.66$\pm$0.78 & 28.60$\pm$0.71 & 28.61$\pm$0.62 & 28.47$\pm$1.05 & 27.62$\pm$0.59 & 4.00 \\
  & & FLDetector & 44.36$\pm$0.59 & 27.71$\pm$4.63 & 44.61$\pm$0.59 & \textbf{44.75$\pm$0.56} & \textbf{44.78$\pm$0.18} & 2.25 \\
  & & \cellcolor{caratrow}\textbf{CARAT ($T=8$)} & \cellcolor{caratrow}44.07$\pm$0.85 & \cellcolor{caratrow}43.09$\pm$0.79 & \cellcolor{caratrow}\textbf{44.80$\pm$0.85} & \cellcolor{caratrow}44.59$\pm$0.99 & \cellcolor{caratrow}43.70$\pm$0.91 & \cellcolor{caratrow}\textbf{2.00} \\
  \bottomrule
  \end{tabular}}
\end{table}
\FloatBarrier

\subsection{Experimental Results}

\subsubsection{Mitigating Conventional Model-Poisoning Attacks}
For $\alpha\in\{1.0,0.5\}$, CARAT has an average rank of $2.50$ under the four conventional attacks. It leads under FangAttack at $\alpha=0.5$ and lies within $0.03$--$0.54$ percentage points of the strongest row in the other seven cells. CARAT also avoids FLDetector's sharp FangAttack degradation and exceeds MultiKrum and FLTrust throughout these two heterogeneity settings. Nevertheless, unprotected Mean has the best average rank at $\alpha=0.5$ and ties FLDetector for the best rank at $\alpha=1.0$, so these results do not establish a uniform improvement over averaging.

At $\alpha=0.1$, CARAT leads MinMax at $44.80\%$, ranks second under FangAttack and MinSum, and reaches an average rank of $2.00$, tied with Mean. Its attack means are within $0.16$--$0.83$ percentage points of the strongest recorded values, with seed-level standard deviations below one percentage point in all four cells. The lower attacker fraction therefore removes the chance-level collapses observed in the archived $20\%$ CARAT runs. This is evidence of stability for the evaluated $10\%$ threat level, not evidence that the same behavior persists as the malicious population grows.

\subsubsection{Robustness to Mimic}
Mimic stresses defenses with repeated benign-looking updates rather than an independently optimized outlier. CARAT reaches $51.27\%$, $49.91\%$, and $43.70\%$ as $\alpha$ decreases from $1.0$ to $0.1$. It ranks second among defenses with recorded Mimic runs in all three settings, trailing FLDetector by $0.46$, $0.92$, and $1.08$ percentage points, respectively, while remaining above NormClipping, MultiKrum, and FLTrust. These results indicate that hidden-task validation retains useful signal against copied benign updates, but FLDetector is consistently stronger in this test.

\subsubsection{Impact of Client Heterogeneity}
Averaged across all five attacks, CARAT obtains $51.30\%$, $49.94\%$, and $44.05\%$ for $\alpha=1.0$, $0.5$, and $0.1$, respectively. The $1.36$-point decline from $\alpha=1.0$ to $0.5$ grows to $5.89$ points from $0.5$ to $0.1$, but the decline is consistent across attacks rather than driven by isolated seed collapses. Client heterogeneity therefore remains an important source of degradation even when the malicious-client fraction is $10\%$.

\subsubsection{Component Contributions}
We conduct a focused component study on CIFAR-100/ResNet18 with $20$ clients, $20\%$ malicious clients, and $\alpha=0.5$. This is a separate, higher-adversary stress test; the available ablation archive does not contain a $10\%$ counterpart. Each variant uses the same five seeds as its within-study control. One task replaces $T=8$ with $T=1$; no clipping disables MAD clipping; no rank penalty sets its coefficient to zero; and no temporal prior removes the temporal regularizer. The ablation covers ALIE, FangAttack, MinMax, and MinSum and was executed independently from the main comparison, so the full CARAT row in Table~\ref{tab:carat-ablation} is the relevant control.

\begin{table}[!htbp]
  \definecolor{caratrow}{gray}{0.88}
  \renewcommand{\arraystretch}{1.05}
  \caption{CARAT component study on CIFAR-100/ResNet18 with $20\%$ malicious clients and non-i.i.d.\ $\alpha=0.5$. Values are final test accuracy (\%) over five seeds. The final column averages the four attack means; higher is better.}
  \label{tab:carat-ablation}
  \setlength{\tabcolsep}{7pt}
  \centering\small
  \begin{tabular}{lccccc}
  \toprule
  \textbf{Variant} & \textbf{ALIE} & \textbf{FangAttack} & \textbf{MinMax} & \textbf{MinSum} & \textbf{Mean} \\
  \midrule
  \cellcolor{caratrow}\textbf{Full CARAT} & \cellcolor{caratrow}\textbf{46.81$\pm$1.06} & \cellcolor{caratrow}\textbf{47.59$\pm$0.61} & \cellcolor{caratrow}48.20$\pm$1.12 & \cellcolor{caratrow}47.82$\pm$0.99 & \cellcolor{caratrow}\textbf{47.61} \\
  One task ($T=1$) & 46.68$\pm$0.93 & 47.12$\pm$0.76 & 48.17$\pm$0.95 & \textbf{48.24$\pm$0.82} & 47.55 \\
  No clipping & 46.26$\pm$0.62 & 46.79$\pm$0.85 & \textbf{48.30$\pm$1.01} & 48.01$\pm$1.00 & 47.34 \\
  No rank penalty & 46.36$\pm$1.11 & 46.96$\pm$0.59 & 48.09$\pm$0.49 & 47.90$\pm$1.00 & 47.33 \\
  No temporal prior & 46.33$\pm$0.80 & 47.42$\pm$1.18 & 48.12$\pm$1.22 & 47.93$\pm$1.04 & 47.45 \\
  \bottomrule
  \end{tabular}
\end{table}

Full CARAT achieves the highest cross-attack mean, but the separation is modest: its advantage over the reduced variants ranges from $0.06$ to $0.28$ percentage points and is smaller than the seed-level variation in individual attack cells. Full CARAT leads under ALIE and FangAttack, whereas no clipping is highest under MinMax and the one-task variant is highest under MinSum. In particular, $T=1$ nearly matches $T=8$ in aggregate. The ablation therefore suggests a small combined benefit in this setting, rather than a large and monotonic contribution from every component.

\subsubsection{Runtime Overhead}
Figure~\ref{fig:verification-curves} reports a separate, earlier CIFAR-100/i.i.d.\ ALIE verification run with three seeds. CARAT records an average per-round time of $9.94$ seconds, compared with $2.80$ seconds for Mean under attack, or approximately $3.6\times$ higher wall-clock cost. This premium is consistent with CARAT evaluating every candidate update on repeated hidden tasks. The diagnostic predates the five-seed main sweep, so it quantifies the implementation tradeoff in that verification setting rather than providing hardware-normalized timing for Table~\ref{tab:cifar100-main-acc}.

\begin{figure*}[!t]
  \centering
  \safeincludegraphics[width=0.98\textwidth]{figures/carat_verification.pdf}
  \caption{Verification-scale CIFAR-100/i.i.d.\ results under ALIE. Left: round-wise test accuracy over three seeds; the dashed clean Mean curve is a single-seed no-attack reference. Right: average per-round wall-clock time with seed-level error bars. This earlier run is an overhead diagnostic and is not part of the main five-seed comparison.}
  \label{fig:verification-curves}
\end{figure*}

\subsubsection{Scope of the Evidence}
The main experiments cover one dataset--model pair, $20$ fully participating clients, a $10\%$ malicious-client ratio, three non-i.i.d.\ concentrations, a known attacker fraction used by CARAT's weight cap, and five untargeted attacks. They do not cover partial participation, larger federations, backdoor or data-poisoning attacks, misspecified attacker fractions, or adaptive attackers that infer the hidden tasks. The missing Mimic--Mean cell, overlap between server-reference and client-training examples, unequal CARAT/FLTrust reference budgets, and the lack of a $10\%$ counterpart to the component study further limit comparison. Together with the strong Mean results on the four conventional attacks and the small ablation gaps, these boundaries restrict the empirical claim to the evaluated threat level rather than universal robustness.

\FloatBarrier

\section{Conclusion}
\label{sec:conclusion}

CARAT validates candidate updates through hidden, class-balanced task certificates rather than relying only on geometric centrality or a single trusted reference direction. In the five-seed CIFAR-10/VGG19 and CIFAR-100/ResNet18 evaluation with $10\%$ malicious clients, CARAT attains an average rank of $1.75$--$2.75$ under the four conventional attacks and ranks first or second among defenses with Mimic results across all six dataset--heterogeneity settings. The CIFAR-100 results remain stable across seeds, while the CIFAR-10 results expose isolated chance-level failures under MinSum and MinMax as heterogeneity increases. Together with the strong unprotected Mean results on the conventional attacks and the small gaps in the separate $20\%$ component study, this evidence supports CARAT as a conditionally competitive reference-assisted mechanism within the evaluated threat level. Disjoint reference data, matched auxiliary-data budgets, adaptive attacks, and systematic evaluation at larger malicious-client fractions are important next steps.



% \bibliographystyle{plainnat}
\bibliography{ref}
\bibliographystyle{iclr2027_conference}

\newpage
\appendix

\paragraph{Appendix organization.}
The appendix is organized to make each claim auditable. Appendix~\ref{sec:appendix-algorithm-theory} gives the round-level algorithm, implementation invariants, and proofs of the structural properties used in the main text. Appendix~\ref{sec:appendix-experiments} gives empirical support for the completed claim: protocol details, CARAT hyperparameters, the separate component study, final-loss results, runtime diagnostics, and reproducibility information.

\section{Algorithmic and Theoretical Details}
\label{sec:appendix-algorithm-theory}

This appendix provides the technical material needed to reproduce and interpret CARAT. It first fixes the round-level notation and server inputs, then gives executable pseudocode, implementation invariants, and structural propositions for the scoring and weighting rule. These results are deliberately stated as invariants and bounds for the proposed aggregation procedure, not as a complete Byzantine-robustness guarantee against arbitrary adaptive attackers.

\subsection{Notation and Round Inputs}
\label{sec:appendix-notation}

All quantities in this appendix refer to one communication round $r$ unless stated otherwise. The server receives $n=|\mathcal{S}_r|$ client messages and converts them into vector updates $\Delta_i^r$. CARAT keeps three versions of each update:
\[
\Delta_i^r
\quad\text{(raw client update)},\qquad
\bar{\Delta}_i^r
\quad\text{(clipped update)},\qquad
\tilde{\Delta}_i^r
\quad\text{(common-radius scoring copy)}.
\]
Only $\bar{\Delta}_i^r$ is used in the final server update. The normalized copy $\tilde{\Delta}_i^r$ is used only to evaluate hidden-task utility, so the validation score is less sensitive to residual norm differences after clipping.

The server maintains a labeled reference set $\mathcal{D}^{\mathrm{ref}}$ and derives from it a working probe pool $\mathcal{P}$. Each hidden task is a small approximately class-balanced subset sampled from $\mathcal{P}$. The raw certificate matrix is $C^{\mathrm{raw}}\in\mathbb{R}^{n\times T}$, where $T$ is the number of hidden tasks and $C^{\mathrm{raw}}_{i,t}=c_{i,t}$ is the loss reduction obtained by candidate update $i$ on hidden task $t$. After robust column-wise standardization, the optimization uses $C\in\mathbb{R}^{n\times T}$. The vector $\gamma\in\mathbb{R}^n_{\ge 0}$ denotes coordinate-rank anomaly penalties. The vector $\alpha^{\mathrm{prev}}$ is the previous round's aggregation weight vector and acts as a temporal prior.

\subsection{CARAT Round Algorithm}
\label{sec:appendix-algorithm}

Algorithm~\ref{alg:carat-round} summarizes one CARAT server aggregation round. The algorithm is server-side: clients do not observe the realized hidden tasks sampled from $\mathcal{P}$.

\begin{algorithm}[t]
\caption{CARAT server aggregation at round $r$}
\label{alg:carat-round}
\begin{algorithmic}[1]
\REQUIRE Current global model $w_r$, client messages $\{u_i^r\}_{i=1}^n$, previous-round weights $\alpha^{\mathrm{prev}}$, server reference set $\mathcal{D}^{\mathrm{ref}}$, probe pool $\mathcal{P}$, CARAT hyperparameters
\ENSURE Aggregated update $\Delta_{\mathrm{agg}}$, updated temporal prior $\alpha^{\mathrm{prev}}$
\IF{$\mathcal{P}$ is uninitialized}
    \STATE Construct a class-balanced probe pool $\mathcal{P}$ from $\mathcal{D}^{\mathrm{ref}}$
\ENDIF
\IF{$r=1$}
    \STATE $\alpha^{\mathrm{prev}} \leftarrow \frac{1}{n}\mathbf{1}$
\ENDIF
\FOR{$i=1$ to $n$}
    \STATE Convert $u_i^r$ into a vector-form delta $\Delta_i^r$
\ENDFOR
\STATE Compute the clipping threshold $\tau_r$ from $\{\|\Delta_i^r\|_2\}_{i=1}^n$
\FOR{$i=1$ to $n$}
    \STATE $\bar{\Delta}_i^r \leftarrow \min\!\left(1,\frac{\tau_r}{\|\Delta_i^r\|_2+\varepsilon}\right)\Delta_i^r$
\ENDFOR
\STATE Choose the evaluation radius $q_r$ from a quantile of positive clipped norms
\FOR{$i=1$ to $n$}
    \STATE $\tilde{\Delta}_i^r \leftarrow q_r \bar{\Delta}_i^r / (\|\bar{\Delta}_i^r\|_2+\varepsilon)$
\ENDFOR
\STATE Sample $T$ hidden tasks from $\mathcal{P}$ using class-balanced sampling
\FOR{$t=1$ to $T$}
    \STATE Evaluate the baseline loss $\ell_t(w_r)$
\ENDFOR
\FOR{$i=1$ to $n$}
    \FOR{$t=1$ to $T$}
        \STATE $c_{i,t} \leftarrow \ell_t(w_r) - \ell_t(w_r + \tilde{\Delta}_i^r)$
    \ENDFOR
\ENDFOR
\STATE Robustly standardize certificate columns to obtain $C$
\STATE Compute rank penalties $\gamma$ from repeated coordinate subsamples of $\{\bar{\Delta}_i^r\}_{i=1}^n$
\STATE $\alpha^{(0)} \leftarrow \alpha^{\mathrm{prev}}$
\FOR{$k=0$ to $K-1$}
    \STATE $z \leftarrow C^\top \alpha^{(k)}$
    \STATE $\pi \leftarrow \softmin$ task weights induced by $z$ and $\beta$
    \STATE $g \leftarrow C\pi - \lambda_r \gamma - 2\lambda_p(\alpha^{(k)}-\alpha^{\mathrm{prev}})$
    \STATE $\alpha^{(k+1)} \leftarrow \Pi_{\mathcal{W}(\hat{\rho})}\!\left(\alpha^{(k)}+\eta g\right)$
\ENDFOR
\STATE $\alpha^\star \leftarrow \alpha^{(K)}$ and $\alpha^{\mathrm{prev}} \leftarrow \alpha^\star$
\STATE $\Delta_{\mathrm{agg}} \leftarrow \sum_{i=1}^n \alpha_i^\star \bar{\Delta}_i^r$
\STATE \textbf{return} $\Delta_{\mathrm{agg}}, \alpha^{\mathrm{prev}}$
\end{algorithmic}
\end{algorithm}

\subsection{Implementation Invariants}
\label{sec:appendix-implementation}

This subsection records implementation choices that affect reproducibility but are too low-level for the main method section. The purpose is to make the server-side scoring rule auditable, not to introduce additional assumptions beyond those stated in the threat model.

\paragraph{Component rationale.}
\label{sec:appendix-component-rationale}
Table~\ref{tab:component-rationale} summarizes why each CARAT component is included. The design is not intended as an arbitrary stack of heuristics: each term targets a distinct failure mode that appears when heterogeneous benign drift and malicious manipulation coexist.

\begin{table*}[!htbp]
\centering
\caption{Design rationale for the CARAT aggregation pipeline. Each component addresses a different ambiguity in non-i.i.d.\ robust aggregation.}
\label{tab:component-rationale}
\renewcommand{\arraystretch}{1.08}
\small
\resizebox{\textwidth}{!}{%
\begin{tabular}{p{0.23\textwidth}p{0.33\textwidth}p{0.36\textwidth}}
\toprule
\textbf{Component} & \textbf{Failure mode addressed} & \textbf{Role in CARAT} \\
\midrule
MAD clipping and common-radius scoring & Large-norm or scale-manipulated updates can dominate validation scores, while benign non-i.i.d.\ clients may naturally have different update norms. & Separates the scoring copy from the final clipped update, making hidden-task evaluation less sensitive to update magnitude. \\
Hidden class-balanced tasks & A single global score or trusted direction can miss class-local degradation and can confuse semantic usefulness with geometric centrality. & Scores each candidate by loss reduction on multiple hidden semantic slices, producing a task-level certificate matrix. \\
Robust task standardization & Hidden tasks can have different baseline loss scales and difficulty levels. & Makes certificate columns comparable while preserving within-task ordering. \\
Coordinate-rank penalty & A client may look useful on part of the reference signal while remaining structurally extreme in update space. & Penalizes consistently outlying coordinate ranks without rewarding unusual centrality. \\
Capped-simplex weights & Even a high-scoring client should not receive unbounded influence in a Byzantine setting. & Limits per-client weight according to an attacker-fraction prior and bounds direct malicious mass. \\
Temporal regularization & Round-to-round weight oscillations can overreact to noisy probe evaluations. & Encourages stable trust allocation unless hidden-task evidence supports a change. \\
\bottomrule
\end{tabular}
}
\end{table*}

\paragraph{Persistent state and initialization.}
CARAT maintains two state variables across rounds: the probe pool $\mathcal{P}$ derived from $\mathcal{D}^{\mathrm{ref}}$, and the previous-round weight vector $\alpha^{\mathrm{prev}}$. At the first round, $\alpha^{\mathrm{prev}}$ is initialized to the uniform vector. This makes the first aggregation unbiased across clients while allowing later rounds to use temporal consistency once weight trajectories have been observed.

\paragraph{Reference-pool construction and hidden-task sampling.}
Let $\mathcal{P}_c \subseteq \mathcal{P}$ denote the subset of probe examples with class label $c$. A hidden task is constructed by sampling a small number of examples from each available class pool and concatenating them into one balanced micro-task. If the desired number of samples per class is not available without replacement, CARAT falls back to replacement and, if needed, to a random subset of the probe pool rather than aborting the round. This fail-soft design preserves the intended semantic diversity of the task sampler while ensuring that the server always obtains a well-defined certificate matrix.

\paragraph{Robust score calibration.}
The raw certificate matrix is transformed column-wise before optimization:
\[
C_{i,t}
=
\frac{C^{\mathrm{raw}}_{i,t}-\operatorname{median}_j(C^{\mathrm{raw}}_{j,t})}
{1.4826 \cdot \operatorname{MAD}_j(C^{\mathrm{raw}}_{j,t})+\varepsilon}.
\]
This removes task-specific scale effects while preserving within-task ordering by Proposition~\ref{prop:task-standardization}. Rank penalties are calibrated similarly. For each coordinate subsample $S_s$, let
\[
\rho_i^{(s)}=
\frac{1}{|S_s|}\sum_{j \in S_s}
\left|
\frac{\operatorname{rank}_{i,j}}{n-1} - \frac{1}{2}
\right|.
\]
CARAT then computes a standardized rank score
\[
z_i^{(s)}=
\frac{\rho_i^{(s)}-\operatorname{median}_\ell(\rho_\ell^{(s)})}
{1.4826 \cdot \operatorname{MAD}_\ell(\rho_\ell^{(s)})+\varepsilon},
\]
and aggregates subsamples through
\[
\gamma_i=\max\left(0,\operatorname{median}_s z_i^{(s)}\right).
\]
Thus a client is penalized only when it is consistently extreme across repeated coordinate views, not because of a single noisy subsample.

\paragraph{Projected optimization.}
CARAT initializes the projected optimizer with $\alpha^{(0)}=\alpha^{\mathrm{prev}}$. If
\[
z(\alpha)=C^\top \alpha,
\qquad
\pi_t(\alpha)=\frac{\exp(-\beta z_t(\alpha))}{\sum_{\ell=1}^T \exp(-\beta z_\ell(\alpha))},
\]
then the ascent direction used in projected updates is
\[
g(\alpha)=C\pi(\alpha)-\lambda_r \gamma-2\lambda_p(\alpha-\alpha^{\mathrm{prev}}).
\]
With optimizer step size $\eta$, one iteration takes the form
\[
\alpha^{(k+1)}=
\Pi_{\mathcal{W}(\hat{\rho})}
\left(\alpha^{(k)}+\eta g(\alpha^{(k)})\right).
\]
The projection maintains feasibility throughout optimization, so every intermediate iterate satisfies nonnegativity, unit-sum, and per-client cap constraints.

\paragraph{Diagnostics.}
The implementation records learned weights, task scores, softmin task weights, rank penalties, clipping thresholds, evaluation radii, and hidden-task identifiers. These diagnostics distinguish three qualitatively different behaviors: benign but diverse clients that retain moderate weight, harmful clients that lose certificate support across hidden tasks, and structurally extreme clients whose influence is reduced mainly through the rank penalty.

\subsection{Proofs and Structural Properties}
\label{sec:appendix-theory}

The following propositions are grouped by the role they play in CARAT. Propositions~\ref{prop:capped-simplex}--\ref{prop:malicious-mass-bound} characterize feasible weighted aggregation and bounded influence after clipping. Propositions~\ref{prop:directional-invariance}--\ref{prop:task-standardization} characterize the hidden-task evaluation geometry. Propositions~\ref{prop:softmin}--\ref{prop:pairwise-transfer} characterize the optimization objective, and Proposition~\ref{prop:complexity} records the per-round cost.

\begin{proposition}[Feasibility and support of the capped simplex]
\label{prop:capped-simplex}
Let
\[
\mathcal{W}(\hat{\rho})=
\left\{
\alpha \in \R^n:
\alpha_i \ge 0,\;
\sum_{i=1}^n \alpha_i = 1,\;
\alpha_i \le \frac{1}{(1-\hat{\rho})n}
\right\},
\]
with $\hat{\rho}<1$. Then $\mathcal{W}(\hat{\rho})$ is non-empty, and every $\alpha \in \mathcal{W}(\hat{\rho})$ satisfies
\[
|\operatorname{supp}(\alpha)| \ge (1-\hat{\rho})n.
\]
\end{proposition}

\begin{proof}
Let $u = 1/((1-\hat{\rho})n)$. Since $\hat{\rho}<1$, we have $u \ge 1/n$, so the uniform weight vector is feasible and the set is non-empty. For the support bound, if only $k$ coordinates of $\alpha$ are positive, then
\[
1=\sum_{i=1}^n \alpha_i \le k u = \frac{k}{(1-\hat{\rho})n},
\]
which implies $k \ge (1-\hat{\rho})n$. Therefore CARAT can suppress suspicious clients, but cannot collapse the aggregate onto an arbitrarily small subset of participants.
\end{proof}

\begin{proposition}[Norm bound for the final aggregate]
\label{prop:aggregate-bound}
Suppose each clipped update satisfies $\|\bar{\Delta}_i\|_2 \le \tau_r$ and $\alpha \in \mathcal{W}(\hat{\rho})$. Then the final aggregate
\[
\Delta_{\mathrm{agg}}=\sum_{i=1}^n \alpha_i \bar{\Delta}_i
\]
satisfies
\[
\|\Delta_{\mathrm{agg}}\|_2 \le \tau_r.
\]
\end{proposition}

\begin{proof}
Because $\alpha_i \ge 0$ and $\sum_i \alpha_i=1$, the aggregate is a convex combination of clipped updates. By the triangle inequality,
\[
\|\Delta_{\mathrm{agg}}\|_2
\le
\sum_{i=1}^n \alpha_i \|\bar{\Delta}_i\|_2
\le
\sum_{i=1}^n \alpha_i \tau_r
=
\tau_r.
\]
Thus the final server step inherits the clipping radius as a hard norm upper bound.
\end{proof}

\begin{proposition}[Influence bound from malicious weight mass]
\label{prop:malicious-mass-bound}
Let $\mathcal{M}$ be the malicious clients selected in a round and let
\[
q_{\mathcal{M}}=\sum_{i\in\mathcal{M}}\alpha_i
\]
be their total CARAT weight. Suppose every clipped update satisfies $\|\bar{\Delta}_i\|_2\le\tau_r$. Then the direct malicious contribution obeys
\[
\left\|\sum_{i\in\mathcal{M}}\alpha_i\bar{\Delta}_i\right\|_2
\le q_{\mathcal{M}}\tau_r.
\]
If $q_{\mathcal{M}}<1$ and
\[
\Delta_{\mathcal{B}}=
\sum_{i\notin\mathcal{M}}\frac{\alpha_i}{1-q_{\mathcal{M}}}\bar{\Delta}_i
\]
denotes the benign-only aggregate obtained by renormalizing the benign weights, then
\[
\|\Delta_{\mathrm{agg}}-\Delta_{\mathcal{B}}\|_2
\le 2q_{\mathcal{M}}\tau_r.
\]
\end{proposition}

\begin{proof}
The first inequality follows from the triangle inequality:
\[
\left\|\sum_{i\in\mathcal{M}}\alpha_i\bar{\Delta}_i\right\|_2
\le
\sum_{i\in\mathcal{M}}\alpha_i\|\bar{\Delta}_i\|_2
\le q_{\mathcal{M}}\tau_r.
\]
For the second inequality, write
\[
\Delta_{\mathrm{agg}}=
(1-q_{\mathcal{M}})\Delta_{\mathcal{B}}
+\sum_{i\in\mathcal{M}}\alpha_i\bar{\Delta}_i.
\]
Therefore
\[
\|\Delta_{\mathrm{agg}}-\Delta_{\mathcal{B}}\|_2
\le
q_{\mathcal{M}}\|\Delta_{\mathcal{B}}\|_2
+\left\|\sum_{i\in\mathcal{M}}\alpha_i\bar{\Delta}_i\right\|_2.
\]
Since $\Delta_{\mathcal{B}}$ is a convex combination of clipped benign updates, $\|\Delta_{\mathcal{B}}\|_2\le\tau_r$ by Proposition~\ref{prop:aggregate-bound}. Combining this with the first inequality gives the stated bound.
\end{proof}

\begin{proposition}[Directional invariance of certificate scoring]
\label{prop:directional-invariance}
Suppose two nonzero clipped updates satisfy $\bar{\Delta}_i = a \bar{\Delta}_j$ for some scalar $a>0$. Then their normalized versions in Eq.~\eqref{eq:normalize} coincide, i.e.,
\[
\tilde{\Delta}_i=\tilde{\Delta}_j.
\]
Consequently, their hidden-task certificates are identical:
\[
c_{i,t}=c_{j,t}\qquad \text{for all } t \in \{1,\dots,T\}.
\]
\end{proposition}

\begin{proof}
Eq.~\eqref{eq:normalize} rescales every nonzero clipped update to the common radius $q_r$. Positive collinearity is therefore removed by normalization:
\[
q_r\frac{\bar{\Delta}_i}{\|\bar{\Delta}_i\|_2}
=
q_r\frac{a\bar{\Delta}_j}{\|a\bar{\Delta}_j\|_2}
=
q_r\frac{\bar{\Delta}_j}{\|\bar{\Delta}_j\|_2}.
\]
Substituting the resulting equality into Eq.~\eqref{eq:cert} gives identical task scores. Thus, after clipping, the certificate stage depends on direction rather than residual norm magnitude.
\end{proof}

\begin{proposition}[Task-wise standardization preserves within-task ordering]
\label{prop:task-standardization}
Fix a task $t$ and define a standardized certificate score
\[
\hat{c}_{i,t}=\frac{c_{i,t}-m_t}{s_t+\varepsilon},
\]
where $m_t \in \R$ and $s_t>0$ depend only on task $t$ and not on client $i$. Then for any clients $i,j$,
\[
c_{i,t} \ge c_{j,t}
\quad \Longleftrightarrow \quad
\hat{c}_{i,t} \ge \hat{c}_{j,t}.
\]
Moreover,
\[
\hat{c}_{i,t}-\hat{c}_{j,t}
=
\frac{c_{i,t}-c_{j,t}}{s_t+\varepsilon}.
\]
\end{proposition}

\begin{proof}
For a fixed task $t$, the map $x \mapsto (x-m_t)/(s_t+\varepsilon)$ is affine with strictly positive slope because $s_t+\varepsilon>0$. Hence it preserves ordering. The difference identity follows by direct subtraction. Therefore column-wise robust standardization calibrates cross-task scale without changing the within-task ranking induced by the raw certificate scores.
\end{proof}

\begin{proposition}[Softmin bounds and gradient form]
\label{prop:softmin}
Let $s \in \R^T$ and define
\[
\softmin_{\beta}(s)
= -\frac{1}{\beta}\log\left(\frac{1}{T}\sum_{t=1}^T e^{-\beta s_t}\right).
\]
Then
\[
\min_t s_t \le \softmin_{\beta}(s) \le \min_t s_t + \frac{\log T}{\beta}.
\]
Moreover,
\[
\lim_{\beta \to \infty}\softmin_{\beta}(s)=\min_t s_t,
\qquad
\lim_{\beta \to 0}\softmin_{\beta}(s)=\frac{1}{T}\sum_{t=1}^T s_t.
\]
If $z(\alpha)=C^\top\alpha$ and
\[
\pi_t(\alpha)=\frac{\exp(-\beta z_t(\alpha))}{\sum_{\ell=1}^T \exp(-\beta z_\ell(\alpha))},
\]
then
\[
\nabla_\alpha \softmin_{\beta}(C^\top\alpha)=C\pi(\alpha).
\]
\end{proposition}

\begin{proof}
Let $m=\min_t s_t$. Since
\[
e^{-\beta m}/T \le \frac{1}{T}\sum_{t=1}^T e^{-\beta s_t} \le e^{-\beta m},
\]
applying $-(1/\beta)\log(\cdot)$ yields
\[
m \le \softmin_{\beta}(s) \le m+\frac{\log T}{\beta}.
\]
The limit $\beta \to \infty$ follows from the vanishing gap bound. For $\beta \to 0$, a first-order expansion of the exponential average gives
\[
\frac{1}{T}\sum_{t=1}^T e^{-\beta s_t}
=
1-\beta \frac{1}{T}\sum_{t=1}^T s_t + o(\beta),
\]
so the softmin converges to the arithmetic mean. Differentiating the softmin expression with respect to $z_t$ yields the Gibbs weights $\pi_t(\alpha)$, and the chain rule gives $C\pi(\alpha)$. Hence the ascent direction is a convex combination of certificate columns, with more weight placed on tasks whose current aggregate utility is smaller.
\end{proof}

\begin{proposition}[Temporal prior and projected-iterate invariants]
\label{prop:temporal-prior}
Define
\[
f(\alpha)=\softmin_{\beta}(C^\top\alpha)-\lambda_r \gamma^\top\alpha-\lambda_p\|\alpha-\alpha^{\mathrm{prev}}\|_2^2.
\]
If $\alpha,\alpha' \in \mathcal{W}(\hat{\rho})$ satisfy $C^\top\alpha=C^\top\alpha'$ and $\gamma^\top\alpha=\gamma^\top\alpha'$, then
\[
f(\alpha)-f(\alpha')
=
-\lambda_p\!\left(\|\alpha-\alpha^{\mathrm{prev}}\|_2^2-\|\alpha'-\alpha^{\mathrm{prev}}\|_2^2\right).
\]
In addition, under projected gradient ascent
\[
\alpha^{(k+1)}=\Pi_{\mathcal{W}(\hat{\rho})}\!\left(\alpha^{(k)}+\eta \nabla f(\alpha^{(k)})\right),
\]
every iterate remains in $\mathcal{W}(\hat{\rho})$.
\end{proposition}

\begin{proof}
The first statement follows by direct substitution: if the certificate and rank terms match, the only remaining difference in objective value is the quadratic prior term. CARAT therefore prefers, among utility-equivalent feasible weights, the solution closer to $\alpha^{\mathrm{prev}}$. The second statement is immediate from the projection operator. Hence the optimizer preserves three round invariants throughout: $\alpha_i \ge 0$, $\sum_i \alpha_i =1$, and $\alpha_i \le 1/((1-\hat{\rho})n)$.
\end{proof}

\begin{proposition}[Pairwise weight-transfer condition]
\label{prop:pairwise-transfer}
Let
\[
f(\alpha)=\softmin_{\beta}(C^\top\alpha)-\lambda_r \gamma^\top\alpha-\lambda_p\|\alpha-\alpha^{\mathrm{prev}}\|_2^2
\]
be the CARAT objective, and let $e_a,e_b$ be coordinate vectors for two clients $a$ and $b$. Suppose $\alpha+\epsilon(e_b-e_a)$ remains feasible for sufficiently small $\epsilon>0$. Define
\[
\pi_t(\alpha)=\frac{\exp(-\beta (C^\top\alpha)_t)}
{\sum_{\ell=1}^T \exp(-\beta (C^\top\alpha)_\ell)}.
\]
Then moving infinitesimal weight from client $a$ to client $b$ increases the objective to first order whenever
\[
(C_b-C_a)^\top \pi(\alpha)
-\lambda_r(\gamma_b-\gamma_a)
-2\lambda_p\!\left[
(\alpha_b-\alpha_b^{\mathrm{prev}})
-(\alpha_a-\alpha_a^{\mathrm{prev}})
\right] > 0,
\]
where $C_i$ denotes row $i$ of the standardized certificate matrix.
\end{proposition}

\begin{proof}
Let $d=e_b-e_a$. A first-order expansion gives
\[
f(\alpha+\epsilon d)-f(\alpha)
=\epsilon d^\top\nabla f(\alpha)+o(\epsilon).
\]
Using Proposition~\ref{prop:softmin}, the gradient is
\[
\nabla f(\alpha)=C\pi(\alpha)-\lambda_r \gamma-2\lambda_p(\alpha-\alpha^{\mathrm{prev}}).
\]
Substituting $d=e_b-e_a$ yields exactly the displayed expression. Hence a client gains weight when its hidden-task utility advantage, measured through the current softmin focus $\pi(\alpha)$, outweighs its rank penalty and temporal-prior cost.
\end{proof}

\begin{proposition}[Per-round complexity]
\label{prop:complexity}
Let $n$ be the number of clients, $T$ the number of hidden tasks, $m$ the number of sampled coordinates for rank scoring, $S$ the number of rank subsamples, and $K$ the number of projected-gradient steps. Ignoring constant factors from batch size and model width, the extra defense-side work of CARAT scales as
\[
\text{probe-model forward passes} + O(Snm + KnT).
\]
\end{proposition}

\begin{proof}
CARAT performs one baseline probe evaluation plus $n$ candidate probe evaluations over the $T$ hidden tasks, then computes $S$ rank summaries over $m$ coordinates, and finally executes $K$ projected optimization steps over $n$ weights and $T$ task scores. In the prototype implementation, the probe-model forward passes dominate wall-clock time, while the rank and optimization stages are lower-order terms controlled by $m$, $S$, and $K$.
\end{proof}

\paragraph{Remark on the reference assumption.}
The object $\mathcal{D}^{\mathrm{ref}}$ denotes the labeled server reference set, while $\mathcal{P}$ is a subsampled working pool derived from it for hidden-task construction. This distinction is part of CARAT's threat model and materially strengthens the setting relative to fully data-free robust aggregation.

\subsection{Reproducibility, Ethics, and Availability}
\label{sec:appendix-reproducibility}

\paragraph{Reproducibility.}
All main reported results use the public CIFAR-10 and CIFAR-100 benchmarks and the fixed training protocol described in Section~\ref{sec:experiments}. The aggregation method is specified in Algorithm~\ref{alg:carat-round}, with implementation defaults listed in Appendix~\ref{sec:appendix-experiments}. The main CARAT settings and component variants each use five seeds, and the result archives retain metrics, completion markers, job metadata, and configuration files so that every reported result can be traced to seed-level runs.

\paragraph{Compute resources.}
The main CARAT sweep and reported public-baseline logs were generated through the repository launch scripts on a Compute Canada SLURM cluster; the component study was run with the same experiment interface on a laboratory GPU server. Each run writes a metadata file recording its configuration, command line, seed, Python environment, and CUDA requirement; the main sweep also records the git commit.

\paragraph{Ethics.}
This work studies defenses against model-poisoning attacks in federated learning. The attack implementations are used to evaluate robustness and do not require access to private user data. The main practical risk is dual use: stronger attack simulations may help adversaries tune poisoning strategies. We mitigate this by framing attacks as benchmark stressors and by releasing defensive evaluation code together with the corresponding protocols.

\paragraph{Data availability.}
CIFAR-10 and CIFAR-100 are public datasets. In the reported implementation, CARAT samples a $512$-example server pool from the corresponding public training set used to construct client partitions; those examples are not removed from client local data. This overlap and the unequal CARAT/FLTrust reference budgets are limitations of the current evaluation. A stricter reproduction should additionally test deterministic, disjoint server splits while reporting the resulting reduction in client training data.

\section{Experimental Protocols and Additional Results}
\label{sec:appendix-experiments}

This appendix records the exact scope, CARAT configuration, supplementary component and runtime studies, loss convention, and provenance underlying Section~\ref{sec:experiments}.

\subsection{Completed Empirical Scope}

The completed comparison covers CIFAR-10/VGG19 and CIFAR-100/ResNet18 with $20$ fully participating clients, Dirichlet non-i.i.d.\ partitions, $10\%$ malicious clients, and five attacks: ALIE, FangAttack, MinMax, MinSum, and Mimic. The paper reports an unattacked Mean reference and attacked results for NormClipping, MultiKrum, FLTrust, FLDetector, and CARAT. For each dataset--model pair, the selected CARAT results contain all combinations of three displayed heterogeneity levels, five attacks, and five seeds, for $150$ CARAT runs in total. The defense baselines and clean Mean references also use five seeds.

\subsection{Reported CIFAR Protocol}

\begin{table*}[!htbp]
\centering
\caption{Detailed protocol for the reported CIFAR-10/VGG19 and CIFAR-100/ResNet18 result sets.}
\label{tab:reported-protocol}
\renewcommand{\arraystretch}{1.08}
\small
\resizebox{\textwidth}{!}{%
\begin{tabular}{p{0.22\textwidth}p{0.30\textwidth}p{0.40\textwidth}}
\toprule
\textbf{Component} & \textbf{Setting} & \textbf{Role or qualification} \\
\midrule
Dataset and model & CIFAR-10 / VGG19; CIFAR-100 / ResNet18 & Complementary ten-class and class-rich benchmarks. \\
Client population & $20$ clients, full participation & Fixed participation in every round. \\
Training schedule & SGD; $200$ rounds; $4$ local epochs; learning rate $0.05$; momentum $0.9$; weight decay $5\times10^{-4}$; batch size $64$ & Common optimization schedule for attacked runs. \\
Data heterogeneity & Dirichlet non-i.i.d.\ $\alpha\in\{1,0.5,0.1\}$ & Rows displayed in the main figures. \\
Malicious clients & $10\%$ & Two of the $20$ clients; this fraction is also supplied to CARAT as $\hat{\rho}$. \\
Attack suite & ALIE, FangAttack, MinMax, MinSum, Mimic & Four conventional poisoning attacks and a benign-update copying attack. \\
Compared methods & Clean reference: No Attack + Mean; attacked defenses: NormClipping, MultiKrum, FLTrust, FLDetector, CARAT & Curves in Figures~\ref{fig:table1-curves-cifar100} and~\ref{fig:table1-curves-cifar10}. \\
Seeds and source policy & Five seeds for every reported result & All attacked defenses and clean references use seeds $42$--$46$. \\
Server examples & CARAT pool: $512$; FLTrust: $100$ & Budgets are unequal; examples are sampled from the training set and not excluded from client partitions. \\
Reported statistic & Round-wise accuracy, mean $\pm$ sample standard deviation & The discussion compares round-200 means; rankings exclude the clean reference. \\
\bottomrule
\end{tabular}
}
\end{table*}

\subsection{CARAT Hyperparameters}

Table~\ref{tab:carat-hyperparameters} lists the CARAT defaults used in the reported result set. The attacker prior is set to the actual malicious-client fraction; robustness to a misspecified fraction is not evaluated here.

\begin{table*}[!htbp]
\centering
\caption{CARAT hyperparameters used in the reported CIFAR-10/VGG19 and CIFAR-100/ResNet18 runs.}
\label{tab:carat-hyperparameters}
\renewcommand{\arraystretch}{1.08}
\small
\resizebox{\textwidth}{!}{%
\begin{tabular}{p{0.24\textwidth}p{0.24\textwidth}p{0.42\textwidth}}
\toprule
\textbf{Component} & \textbf{Default value} & \textbf{Purpose} \\
\midrule
Reference pool & $512$ training-set examples & Pool from which hidden tasks are sampled; it is not disjoint from client partitions. \\
Hidden-task evaluation & $8$ tasks, $4$ examples per class, batch size $128$ & Builds repeated class-balanced certificates; each task contains $40$ examples on CIFAR-10 and $400$ on CIFAR-100. \\
Normalization and clipping & Probe-radius quantile $0.5$; MAD clipping factor $2.5$ & Places candidates on a common evaluation radius and limits residual norm influence. \\
Feasible weight set & Attacker prior $\hat{\rho}=0.1$ & Caps per-client weight using the known experimental attacker fraction. \\
Task and anomaly weighting & Temperature $\beta=12.0$; rank weight $0.20$; prior weight $0.05$ & Balances weakest-task utility, rank penalty, and temporal smoothness. \\
Rank diagnostics & $4096$ sampled coordinates, $5$ rank subsamples & Estimates persistent coordinate-wise extremeness. \\
Projected optimization & $60$ projected steps, step size $0.5$ & Solves the capped-simplex weighting problem. \\
\bottomrule
\end{tabular}
}
\end{table*}

\subsection{Separate Component Study}
\label{sec:appendix-component-study}

We conduct a focused component study on CIFAR-100/ResNet18 with $20$ clients, $20\%$ malicious clients, and $\alpha=0.5$. This is a separate, higher-adversary stress test; the available archive does not contain a $10\%$ counterpart. Each variant uses the same five seeds as its within-study control. One task replaces $T=8$ with $T=1$; no clipping disables MAD clipping; no rank penalty sets its coefficient to zero; and no temporal prior removes the temporal regularizer. The study covers ALIE, FangAttack, MinMax, and MinSum and was executed independently from the main comparison.

\begin{table}[!htbp]
  \definecolor{caratrow}{gray}{0.88}
  \renewcommand{\arraystretch}{1.05}
  \caption{CARAT component study on CIFAR-100/ResNet18 with $20\%$ malicious clients and non-i.i.d.\ $\alpha=0.5$. Values are final test accuracy (\%) over five seeds. The final column averages the four attack means; higher is better.}
  \label{tab:carat-ablation}
  \setlength{\tabcolsep}{7pt}
  \centering\small
  \begin{tabular}{lccccc}
  \toprule
  \textbf{Variant} & \textbf{ALIE} & \textbf{FangAttack} & \textbf{MinMax} & \textbf{MinSum} & \textbf{Mean} \\
  \midrule
  \cellcolor{caratrow}\textbf{Full CARAT} & \cellcolor{caratrow}\textbf{46.81$\pm$1.06} & \cellcolor{caratrow}\textbf{47.59$\pm$0.61} & \cellcolor{caratrow}48.20$\pm$1.12 & \cellcolor{caratrow}47.82$\pm$0.99 & \cellcolor{caratrow}\textbf{47.61} \\
  One task ($T=1$) & 46.68$\pm$0.93 & 47.12$\pm$0.76 & 48.17$\pm$0.95 & \textbf{48.24$\pm$0.82} & 47.55 \\
  No clipping & 46.26$\pm$0.62 & 46.79$\pm$0.85 & \textbf{48.30$\pm$1.01} & 48.01$\pm$1.00 & 47.34 \\
  No rank penalty & 46.36$\pm$1.11 & 46.96$\pm$0.59 & 48.09$\pm$0.49 & 47.90$\pm$1.00 & 47.33 \\
  No temporal prior & 46.33$\pm$0.80 & 47.42$\pm$1.18 & 48.12$\pm$1.22 & 47.93$\pm$1.04 & 47.45 \\
  \bottomrule
  \end{tabular}
\end{table}

Full CARAT has the highest cross-attack mean, but its advantage over the reduced variants is only $0.06$--$0.28$ percentage points and is smaller than the seed-level variation in individual attack settings. Full CARAT leads under ALIE and FangAttack; no clipping leads under MinMax; and one task leads under MinSum. Thus, $T=1$ nearly matches $T=8$, and the study supports only a small combined benefit in this higher-adversary setting.

\FloatBarrier
\subsection{Runtime Diagnostic}
\label{sec:appendix-runtime}

Figure~\ref{fig:verification-curves} reports a separate, earlier CIFAR-100/i.i.d.\ ALIE verification run with three seeds. CARAT records an average per-round time of $9.94$ seconds, compared with $2.80$ seconds for Mean under attack, or approximately $3.6\times$ higher wall-clock cost. This premium is consistent with evaluating every candidate update on repeated hidden tasks. The diagnostic predates the five-seed sweep and is not hardware-normalized for the main comparison.

\begin{figure*}[!b]
  \centering
  \safeincludegraphics[width=0.98\textwidth]{figures/carat_verification.pdf}
  \caption{Verification-scale CIFAR-100/i.i.d.\ results under ALIE. Left: round-wise test accuracy over three seeds; the dashed clean Mean curve is a single-seed no-attack reference. Right: average per-round wall-clock time with seed-level error bars. This earlier run is an overhead diagnostic and is not part of the main five-seed comparison.}
  \label{fig:verification-curves}
\end{figure*}
\FloatBarrier

\subsection{CARAT Final-Loss Results}

The new exporter records evaluation loss as the raw mean cross-entropy. Earlier baseline logs stored cross-entropy with an additional batch-size normalization, so their numerical loss values are not comparable. We therefore avoid a cross-defense loss ranking and report only the internally consistent five-seed CARAT values in Table~\ref{tab:carat-main-loss}. Accuracy in Figures~\ref{fig:table1-curves-cifar100} and~\ref{fig:table1-curves-cifar10} remains the cross-defense metric.

\begin{table}[!htbp]
\centering
\caption{Five-seed CARAT final-round raw mean cross-entropy with $10\%$ malicious clients. Lower is better. A dash marks a cell containing at least one non-finite seed loss.}
\label{tab:carat-main-loss}
\renewcommand{\arraystretch}{1.08}
\small
\setlength{\tabcolsep}{4pt}
\resizebox{\textwidth}{!}{%
\begin{tabular}{ccccccc}
\toprule
\textbf{Dataset / Model} & \textbf{Dirichlet $\alpha$} & \textbf{ALIE} & \textbf{FangAttack} & \textbf{MinMax} & \textbf{MinSum} & \textbf{Mimic} \\
\midrule
\multirow{3}{*}{CIFAR-100 / ResNet18} & $1.0$ & $1.7628 \pm 0.0270$ & $1.7942 \pm 0.0236$ & $1.7572 \pm 0.0323$ & $1.7656 \pm 0.0235$ & $1.7775 \pm 0.0199$ \\
& $0.5$ & $1.8212 \pm 0.0350$ & $1.8400 \pm 0.0308$ & $1.8259 \pm 0.0221$ & $1.8241 \pm 0.0209$ & $1.8252 \pm 0.0327$ \\
& $0.1$ & $2.0843 \pm 0.0381$ & $2.1249 \pm 0.0353$ & $2.0553 \pm 0.0337$ & $2.0643 \pm 0.0415$ & $2.0896 \pm 0.0278$ \\
\midrule
\multirow{3}{*}{CIFAR-10 / VGG19} & $1.0$ & $0.6002 \pm 0.0043$ & $0.6139 \pm 0.0105$ & $0.5741 \pm 0.0150$ & $0.5773 \pm 0.0137$ & $0.6009 \pm 0.0106$ \\
& $0.5$ & $0.6741 \pm 0.0133$ & $0.6960 \pm 0.0088$ & $0.6621 \pm 0.0187$ & -- & $0.6658 \pm 0.0209$ \\
& $0.1$ & $1.2320 \pm 0.1048$ & $1.2611 \pm 0.1128$ & -- & -- & $1.2582 \pm 0.0987$ \\
\bottomrule
\end{tabular}}
\end{table}

All $75$ selected CIFAR-100 runs record finite final-round cross-entropy. On CIFAR-10, the $\alpha=0.5$ MinSum cell and the $\alpha=0.1$ MinMax and MinSum cells each contain one non-finite loss paired with a chance-level $10\%$ accuracy. We mark these cells unavailable rather than report a finite-only mean.

\subsection{Result Provenance and Integrity Checks}

The CARAT client-$20$ sweep contains $960$ completed runs across two datasets, four data-distribution settings, four malicious-client ratios, six attacks, and five seeds. The main paper selects $150$ runs: the $75$ CIFAR-10/VGG19 and $75$ CIFAR-100/ResNet18 runs with $\alpha\in\{1,0.5,0.1\}$, malicious ratio $0.1$, and the five displayed attacks. Every selected metric file contains all $200$ training rounds. Job metadata records code commit \texttt{70920a89}, the launch command, configuration, seed, Python environment, and CUDA requirement. The extracted CARAT-only archive has SHA-256 checksum \texttt{382bac20d9a44dffa41e8b36de1c05943c69d20d4cb1107e6e2963a6bca97134}.

The component study contains $100$ completed runs at malicious ratio $0.2$: five variants, four attacks, and five seeds, each with $200$ rounds. It was executed independently on the laboratory server and is reported as a higher-adversary stress test rather than as a ratio-matched ablation of the main comparison. Because concurrent jobs make wall-clock measurements sensitive to contention, these runs support the accuracy ablation in Table~\ref{tab:carat-ablation} but are not used for a runtime claim.

\FloatBarrier



\subsection*{AI use statement}

(This section is \textbf{required} and does not count toward the page limit.)

In this work, we used generative AI tools for [tasks with required disclosure].
We have not used generative AI tools for [other tasks with required disclosure],
and [the rest of the required disclosure tasks] are not applicable to this work.
Additionally, we used generative AI tools for [tasks with recommended
disclosure]. We have reviewed all AI-assisted work. [Elaborate. For example, “we
checked LLM-generated research ideas for potential plagiarism through a manual
literature survey”, “LLM-generated code was verified and tested for correctness
by 2 authors”, etc.]. We take responsibility for the final content of this work,
including text, claims or artifacts produced with the aid of generative AI.

See the ICLR 2027 AI Policy for Authors for more details. This statement should
not be more than 1 page.

\subsection*{Ethics statement}

(This section is \textbf{recommended} and does not count toward the page limit.)

If authors feel that their paper submission raises questions regarding the Code
of Ethics, they are encouraged to include a paragraph of Ethics Statement (at
the end of the main text before references) to address potential concerns where
appropriate. Topics include, but are not limited to, studies that involve human
subjects, practices to data set releases, potentially harmful insights,
methodologies and applications, potential conflicts of interest and sponsorship,
discrimination/bias/fairness concerns, privacy and security issues, legal
compliance, and research integrity issues (e.g., IRB, documentation, research
ethics). This statement should not be more than 1 page.

\subsection*{Reproducibility statement}

(This section is \textbf{recommended} and does not count toward the page limit.)

It is important that the work published in ICLR is reproducible. Authors are
strongly encouraged to include a paragraph-long Reproducibility Statement at the
end of the main text (before references) to discuss the efforts that have been
made to ensure reproducibility. This paragraph should not itself describe
details needed for reproducing the results, but rather reference the parts of
the main paper, appendix, and supplemental materials that will help with
reproducibility. For example, for novel models or algorithms, a link to an
anonymous downloadable source code can be submitted as supplementary materials;
for theoretical results, clear explanations of any assumptions and a complete
proof of the claims can be included in the appendix; for any datasets used in
the experiments, a complete description of the data processing steps can be
provided in the supplementary materials. Each of the above are examples of
things that can be referenced in the reproducibility statement.



\subsubsection*{Acknowledgments}
Use unnumbered third level headings for the acknowledgments. All
acknowledgments, including those to funding agencies, go at the end of the paper.






\end{document}
